This paper identifies design parameters that can lead to inaccurate estimates of small effects. Low statistical power not only makes such effects difficult to detect but resulting significant estimates necessarily exaggerate effect sizes on average. Through the literature on short-term health effects of air pollution, we explore the prevalence and policy implications of this issue. Large sample sizes are not sufficient to ensure adequate power and prevent exaggeration. Real-data simulations highlight key additional drivers: the number of exogenous shocks, instrument strength, and outcome count. We propose a workflow to evaluate and mitigate exaggeration risk in non-experimental studies.